\section{September 15, 2026}

We begin with several examples illustrating the difference between
measure-theoretic size and topological size. Throughout, $m$ denotes
Lebesgue measure and $\mathcal L$ denotes the Lebesgue $\sigma$-algebra.

\subsection{Measure and Baire category}

Enumerate the rationals as $\bb Q=\{q_1,q_2,\ldots\}$. For
$n\in\bb N$, define
\[
  U_n=\bigcup_{k=1}^{\infty}
  \left(q_k-2^{-n-k-2},q_k+2^{-n-k-2}\right).
\]
Each $U_n$ is open and dense because it contains $\bb Q$, but
\[
  m(U_n)
  \leq\sum_{k=1}^{\infty}2^{-n-k-1}
  =2^{-n-1}.
\]
The set
\[
  U=\bigcap_{n=1}^{\infty}U_n
\]
is a dense $G_\delta$ set by the Baire category theorem, yet
$m(U)=0$. It is necessarily uncountable: a countable subset of
$\bb R$ is meager, whereas a dense $G_\delta$ subset of $\bb R$ is
comeager.

This example has a complementary form on $[0,1]$. Let
\[
  F_n=[0,1]\setminus U_n,
  \qquad
  F=\bigcup_{n=1}^{\infty}F_n=[0,1]\setminus U.
\]
Because $U_n$ is open and dense, $F_n$ is closed and nowhere dense.
Thus $F$ is meager, being a countable union of nowhere dense sets, but
\[
  m(F)=1.
\]
Thus the topological and measure-theoretic notions of size need not agree.

\subsection{The Cantor set}

The middle-thirds Cantor set $C\subseteq[0,1]$ can be described in two
equivalent ways. First, set $C_0=[0,1]$ and obtain $C_n$ from
$C_{n-1}$ by deleting the open middle third of each interval component.
Then
\[
  C=\bigcap_{n=0}^{\infty}C_n.
\]
Equivalently,
\[
  C=\left\{
    \sum_{j=1}^{\infty}a_j3^{-j}:a_j\in\{0,2\}
  \right\}.
\]

\begin{proposition}\label{prop:cantor-set-properties-sep-15}
  The Cantor set has the following properties.
  \begin{enumerate}[label=(\alph*)]
    \item $m(C)=0$.
    \item $C$ is compact, has empty interior, and is totally
      disconnected.
    \item $|C|=|\bb R|$.
  \end{enumerate}
\end{proposition}

\begin{proof}
  The set $C_n$ consists of $2^n$ closed intervals, each of length
  $3^{-n}$. By continuity from above,
  \[
    m(C)=\lim_{n\to\infty}m(C_n)
    =\lim_{n\to\infty}\left(\frac23\right)^n=0.
  \]
  Each $C_n$ is compact, so their nested intersection $C$ is compact.
  Since every nonempty interval has positive measure, the equality
  $m(C)=0$ implies that $C$ has empty interior. Every connected subset
  of $\bb R$ is an interval; since $C$ contains no nontrivial interval,
  each connected component of $C$ is a singleton. Thus $C$ is totally
  disconnected.

  Finally, the map
  \[
    (b_j)_{j=1}^{\infty}
    \longmapsto
    \sum_{j=1}^{\infty}2b_j3^{-j},
    \qquad b_j\in\{0,1\},
  \]
  is a bijection from $\{0,1\}^{\bb N}$ onto $C$. Since the set of
  binary sequences has the cardinality of $\bb R$, so does $C$.
\end{proof}

The Cantor set supports a function that is continuous and increasing but
whose derivative vanishes almost everywhere.

\begin{definition}[Cantor function]
  If $x\in C$ has ternary expansion
  \[
    x=\sum_{j=1}^{\infty}2b_j3^{-j},
    \qquad b_j\in\{0,1\},
  \]
  define
  \[
    \Phi(x)=\sum_{j=1}^{\infty}b_j2^{-j}.
  \]
  Every component of $[0,1]\setminus C$ is an open interval whose
  endpoints lie in $C$ and have the same value under the preceding
  formula. Extend $\Phi$ to be constant on each such interval. The
  resulting map $\Phi\colon[0,1]\to[0,1]$ is the \emph{Cantor function}.
\end{definition}

The Cantor function is continuous and increasing, with
$\Phi(0)=0$ and $\Phi(1)=1$. On the open set $[0,1]\setminus C$ it is
locally constant. Since $m(C)=0$, it follows that
\[
  \Phi'(x)=0\quad\text{for almost every }x\in[0,1].
\]
Nevertheless,
\[
  \Phi(1)-\Phi(0)=1
  \neq0=\int_0^1\Phi'(x)\,dx.
\]
Thus continuity alone is not sufficient for the strongest form of the
fundamental theorem of calculus.

\begin{figure}[htbp]
  \centering
  \begin{tikzpicture}[x=9cm,y=5cm]
    \draw[->] (-0.025,0) -- (1.07,0) node[right] {$x$};
    \draw[->] (0,-0.04) -- (0,1.08) node[above] {$\Phi(x)$};

    \draw[gray!35,dashed] (0,0.5) -- (1,0.5);
    \draw[gray!35,dashed] (0.333333,0) -- (0.333333,0.5);
    \draw[gray!35,dashed] (0.666667,0) -- (0.666667,0.5);

    \draw[very thick,blue!70!black]
      plot coordinates {
        (0,0)
        (0.037037,0.125) (0.074074,0.125)
        (0.111111,0.25) (0.222222,0.25)
        (0.259259,0.375) (0.296296,0.375)
        (0.333333,0.5) (0.666667,0.5)
        (0.703704,0.625) (0.740741,0.625)
        (0.777778,0.75) (0.888889,0.75)
        (0.925926,0.875) (0.962963,0.875)
        (1,1)
      };

    \foreach \x/\lab in {
      0/0,
      0.333333/{\frac13},
      0.666667/{\frac23},
      1/1
    } {
      \draw (\x,0.012) -- (\x,-0.012)
        node[below=2pt] {$\lab$};
    }
    \foreach \y/\lab in {0.5/{\frac12},1/1} {
      \draw (0.008,\y) -- (-0.008,\y)
        node[left=2pt] {$\lab$};
    }
    \node[blue!70!black,above left] at (0.96,0.9) {$\Phi_3$};
  \end{tikzpicture}
  \caption{The third-stage polygonal approximation $\Phi_3$ to the
    Cantor function. At every subsequent stage, each sloping segment is
    replaced by a smaller copy of the first-stage pattern; the limit is
    constant on every component of $[0,1]\setminus C$.}
  \label{fig:cantor-function-sep-15}
\end{figure}

\begin{example}[A fat Cantor set]
  A fixed middle-fifth construction, whose base-$5$ digits are restricted
  to $\{0,1,3,4\}$, still has measure zero because its $n$th stage has
  measure $(4/5)^n$. To obtain a Cantor-type set of positive measure,
  the deleted proportions must shrink.

  Set $K_0=[0,1]$. At stage $n\geq1$, remove from the center of each
  interval component of $K_{n-1}$ an open interval whose length is
  $5^{-n}$ times the length of that component, and call the remaining
  set $K_n$. The first step is
  \[
    K_1=\left[0,\frac25\right]\cup
        \left[\frac35,1\right].
  \]
  For the limit set $K=\bigcap_nK_n$,
  \[
    m(K)=\prod_{n=1}^{\infty}(1-5^{-n})>0,
  \]
  because $\sum_n5^{-n}<\infty$. On the other hand, the lengths of the
  interval components tend to zero. Hence $K$ is compact, has empty
  interior, and is totally disconnected. It is also perfect and has
  cardinality $|\bb R|$. Such a set is called a \emph{fat Cantor set}.
\end{example}

\subsection{The idea of Lebesgue integration}

Let $(X,\mathcal M,\mu)$ be a measure space. For $E\in\mathcal M$, the
\emph{characteristic function} of $E$ is
\[
  \mathbf 1_E(x)=
  \begin{cases}
    1,&x\in E,\\
    0,&x\notin E.
  \end{cases}
\]
The integral is required to satisfy
\[
  \int_X\mathbf 1_E\,d\mu=\mu(E).
\]

This suggests how to integrate a more general function. Rather than
partitioning the domain into short intervals, as in Riemann integration,
Lebesgue integration partitions the range. For a mesh size
$\varepsilon>0$ and $k\in\bb Z$, consider the level sets
\[
  E_k=\{x\in X:k\varepsilon\leq f(x)<(k+1)\varepsilon\}.
\]
When these sets belong to $\mathcal M$, the function can be approximated
by a simple function, and formally
\[
  \int_X f\,d\mu
  \approx\sum_{k\in\bb Z}k\varepsilon\,\mu(E_k).
\]
The requirement that the level sets be measurable leads to the notion
of a measurable function.

\subsection{Measurable functions}

\begin{definition}
  Let $(X,\mathcal M)$ and $(Y,\mathcal N)$ be measurable spaces. A map
  $f\colon X\to Y$ is \emph{$(\mathcal M,\mathcal N)$-measurable} if
  \[
    f^{-1}(E)\in\mathcal M
  \]
  for every $E\in\mathcal N$.
\end{definition}

For a real- or complex-valued function, the codomain is equipped with
its Borel $\sigma$-algebra. Thus $f\colon X\to\bb R$ is
$\mathcal M$-measurable when it is
$(\mathcal M,\mathcal B_{\bb R})$-measurable, and similarly for
$f\colon X\to\bb C$.

The extended real line
\[
  \overline{\bb R}=\bb R\cup\{-\infty,+\infty\}
\]
is equipped with its Borel $\sigma$-algebra. Its topology may be
described by the metric
\[
  d(x,y)=|\arctan x-\arctan y|,
  \qquad
  \arctan(\pm\infty)=\pm\frac{\pi}{2}.
\]

For functions on $\bb R$, the terminology is worth distinguishing:
\begin{itemize}
  \item $f$ is \emph{Borel measurable} if it is
    $(\mathcal B_{\bb R},\mathcal B_{\bb R})$-measurable;
  \item $f$ is \emph{Lebesgue measurable} if it is
    $(\mathcal L,\mathcal B_{\bb R})$-measurable.
\end{itemize}
The codomain still carries the Borel $\sigma$-algebra in the second
definition. In particular, although a composition of appropriately
measurable maps is measurable, the composition of two Lebesgue
measurable functions $\bb R\to\bb R$ need not be Lebesgue measurable:
the inverse image under the outer function is only known to be a
Lebesgue set, while the inner function is only required to pull back
Borel sets.

\begin{lemma}[Checking a generating family]
  \label{lem:measurability-generators-sep-15}
  Suppose $\mathcal N=\mathcal M(\mathcal E)$ for some
  $\mathcal E\subseteq2^Y$. A map $f\colon X\to Y$ is
  $(\mathcal M,\mathcal N)$-measurable if and only if
  \[
    f^{-1}(E)\in\mathcal M
  \]
  for every $E\in\mathcal E$.
\end{lemma}

\begin{proof}
  Only the reverse implication requires proof. Set
  \[
    \mathcal A
    =\{E\subseteq Y:f^{-1}(E)\in\mathcal M\}.
  \]
  Since preimages preserve complements and countable unions,
  $\mathcal A$ is a $\sigma$-algebra. By hypothesis it contains
  $\mathcal E$, so it contains
  $\mathcal M(\mathcal E)=\mathcal N$.
\end{proof}

\begin{corollary}
  If $X$ and $Y$ are topological spaces and
  $f\colon X\to Y$ is continuous, then $f$ is
  $(\mathcal B_X,\mathcal B_Y)$-measurable.
\end{corollary}

\begin{proof}
  The Borel $\sigma$-algebra $\mathcal B_Y$ is generated by the open
  subsets of $Y$, and the inverse image of an open set under a
  continuous map is open.
\end{proof}

\begin{proposition}[Half-line criterion]
  \label{prop:half-line-measurability-sep-15}
  Let $f\colon X\to\overline{\bb R}$. The following are equivalent.
  \begin{enumerate}[label=(\roman*)]
    \item $f$ is measurable.
    \item $\{x:f(x)>a\}\in\mathcal M$ for every $a\in\bb R$.
    \item $\{x:f(x)\geq a\}\in\mathcal M$ for every $a\in\bb R$.
    \item $\{x:f(x)<a\}\in\mathcal M$ for every $a\in\bb R$.
    \item $\{x:f(x)\leq a\}\in\mathcal M$ for every $a\in\bb R$.
  \end{enumerate}
\end{proposition}

\begin{proof}
  Each of the four corresponding families of half-lines generates the
  Borel $\sigma$-algebra of $\overline{\bb R}$. The conclusion follows
  from Lemma~\ref{lem:measurability-generators-sep-15}.
\end{proof}

Finally, if $(X_i,\mathcal M_i)_{i\in I}$ are measurable spaces, the
product $\sigma$-algebra $\bigotimes_{i\in I}\mathcal M_i$ is precisely
the smallest $\sigma$-algebra on $\prod_iX_i$ for which every coordinate
projection
\[
  \pi_j\colon\prod_{i\in I}X_i\longrightarrow X_j
\]
is measurable. This is the measurable-space analogue of the fact that
the product topology is the smallest topology making all coordinate
projections continuous.
